\chapter{Communication Complexity — Subprotocol}
%%% generated by the blueprint pipeline from TCSlib.CommunicationComplexity.DeterministicCC.Subprotocol
\section{Overview}

This module develops the theory of subprotocol embeddings for deterministic communication
protocols.  It introduces an inductive predicate \texttt{IsSubprotocol} asserting that one
protocol tree is a rooted subtree of another, proves that every protocol with more than one
leaf contains a balanced subprotocol (with a $\frac{1}{3}$--$\frac{2}{3}$ fraction of the
leaves), and provides operations for erasing, pruning, and routing inputs through
subprotocols, culminating in the two-bit \texttt{testSubprotocol} construction.

%%%%%%%%%%%%%%%%%%%%%%%%%%%%%%%%%%%%%%%%%%%%%%%%%%%%%%%%%%%%%%%%%%%%%%%%%%%%%%%
\section{Declarations}

\begin{definition}[Subprotocol predicate]
\lean{CommunicationComplexity.Deterministic.Protocol.IsSubprotocol}
\label{CommunicationComplexity.Deterministic.Protocol.IsSubprotocol}
\leanok
\uses{CommunicationComplexity.Deterministic.Protocol}
The inductive proposition $\texttt{IsSubprotocol}\; s\; p$ asserts that protocol $s$ is a
rooted subtree of protocol $p$.  It is generated by: reflexivity ($s$ is a subtree of
itself), and three step cases recording that $s$ is a subtree of one of the two children of
an Alice-node or the false-child of a Bob-node, and hence a subtree of the parent.
\end{definition}

\begin{lemma}[Transitivity of the subprotocol relation]
\lean{CommunicationComplexity.Deterministic.Protocol.IsSubprotocol.trans}
\label{CommunicationComplexity.Deterministic.Protocol.IsSubprotocol.trans}
\leanok
\uses{CommunicationComplexity.Deterministic.Protocol, CommunicationComplexity.Deterministic.Protocol.IsSubprotocol}
\difficulty{3}
Let $s$, $t$, and $u$ be deterministic two-party communication protocols over the same
input types and output type. If $s$ is a subprotocol of $t$ and $t$ is a subprotocol of
$u$, then $s$ is a subprotocol of $u$. In other words, the subprotocol relation is
transitive.
\end{lemma}

\begin{lemma}[A near-balanced subprotocol]
\lean{CommunicationComplexity.Deterministic.Protocol.balanced_aux}
\proofsource{raoyehudayoff-ch01-deterministic-protocols}{
  pdf-fa052d8c95c0-p004-b003,
  pdf-fa052d8c95c0-p004-b004
}
\label{CommunicationComplexity.Deterministic.Protocol.balanced_aux}
\leanok
\proofstep
  {CommunicationComplexity.Deterministic.Protocol}
  {direct}
  {raoyehudayoff-ch01-deterministic-protocols}
  {pdf-fa052d8c95c0-p003-b001}
\proofstep
  {CommunicationComplexity.Deterministic.Protocol.run}
  {direct}
  {raoyehudayoff-ch01-deterministic-protocols}
  {pdf-fa052d8c95c0-p003-b002}
\proofstep
  {CommunicationComplexity.Deterministic.Protocol.complexity}
  {direct}
  {raoyehudayoff-ch01-deterministic-protocols}
  {pdf-fa052d8c95c0-p003-b003}
\proofstep
  {CommunicationComplexity.Deterministic.Protocol.swap}
  {context}
  {raoyehudayoff-ch01-deterministic-protocols}
  {pdf-fa052d8c95c0-p003-b001}
\proofstep
  {CommunicationComplexity.Deterministic.Protocol.swap_run}
  {context}
  {raoyehudayoff-ch01-deterministic-protocols}
  {pdf-fa052d8c95c0-p003-b002}
\proofstep
  {CommunicationComplexity.Deterministic.Protocol.swap_complexity}
  {context}
  {raoyehudayoff-ch01-deterministic-protocols}
  {pdf-fa052d8c95c0-p003-b003}
\proofstep
  {CommunicationComplexity.Deterministic.Protocol.swap_swap}
  {context}
  {raoyehudayoff-ch01-deterministic-protocols}
  {pdf-fa052d8c95c0-p003-b001}
\proofstep
  {CommunicationComplexity.Deterministic.Protocol.comap}
  {context}
  {raoyehudayoff-ch01-deterministic-protocols}
  {pdf-fa052d8c95c0-p003-b001}
\proofstep
  {CommunicationComplexity.Deterministic.Protocol.comap_run}
  {context}
  {raoyehudayoff-ch01-deterministic-protocols}
  {pdf-fa052d8c95c0-p003-b002}
\proofstep
  {CommunicationComplexity.Deterministic.Protocol.comap_complexity}
  {context}
  {raoyehudayoff-ch01-deterministic-protocols}
  {pdf-fa052d8c95c0-p003-b003}
\proofstep
  {CommunicationComplexity.Deterministic.Protocol.shape}
  {context}
  {raoyehudayoff-ch01-deterministic-protocols}
  {pdf-fa052d8c95c0-p003-b001}
\proofstep
  {CommunicationComplexity.Deterministic.Protocol.numLeaves}
  {context}
  {raoyehudayoff-ch01-deterministic-protocols}
  {pdf-fa052d8c95c0-p003-b005,
    pdf-fa052d8c95c0-p004-b002}
\proofstep
  {CommunicationComplexity.Deterministic.Protocol.IsSubprotocol}
  {context}
  {raoyehudayoff-ch01-deterministic-protocols}
  {pdf-fa052d8c95c0-p004-b003,
    pdf-fa052d8c95c0-p005-b002}
\proofstep
  {CommunicationComplexity.Deterministic.Protocol.IsSubprotocol.trans}
  {context}
  {raoyehudayoff-ch01-deterministic-protocols}
  {pdf-fa052d8c95c0-p004-b003}
\proofstep
  {CommunicationComplexity.Deterministic.Protocol.SubprotocolPath}
  {context}
  {raoyehudayoff-ch01-deterministic-protocols}
  {pdf-fa052d8c95c0-p005-b002}
\proofstep
  {CommunicationComplexity.Deterministic.Protocol.path_exists_of_isSubprotocol}
  {context}
  {raoyehudayoff-ch01-deterministic-protocols}
  {pdf-fa052d8c95c0-p005-b002}
\proofstep
  {CommunicationComplexity.Deterministic.Protocol.choosePath}
  {context}
  {raoyehudayoff-ch01-deterministic-protocols}
  {pdf-fa052d8c95c0-p005-b002}
\proofstep
  {CommunicationComplexity.Deterministic.Protocol.reachXPath}
  {context}
  {raoyehudayoff-ch01-deterministic-protocols}
  {pdf-fa052d8c95c0-p005-b002,
    pdf-fa052d8c95c0-p005-b003}
\proofstep
  {CommunicationComplexity.Deterministic.Protocol.reachYPath}
  {context}
  {raoyehudayoff-ch01-deterministic-protocols}
  {pdf-fa052d8c95c0-p005-b002,
    pdf-fa052d8c95c0-p005-b003}
\proofstep
  {CommunicationComplexity.Deterministic.Protocol.reachX}
  {context}
  {raoyehudayoff-ch01-deterministic-protocols}
  {pdf-fa052d8c95c0-p005-b002,
    pdf-fa052d8c95c0-p005-b003}
\proofstep
  {CommunicationComplexity.Deterministic.Protocol.reachY}
  {context}
  {raoyehudayoff-ch01-deterministic-protocols}
  {pdf-fa052d8c95c0-p005-b002,
    pdf-fa052d8c95c0-p005-b003}
\proofstep
  {CommunicationComplexity.Deterministic.Protocol.testSubprotocol}
  {context}
  {raoyehudayoff-ch01-deterministic-protocols}
  {pdf-fa052d8c95c0-p004-b005}
\proofstep
  {CommunicationComplexity.Deterministic.Protocol.testSubprotocol_complexity}
  {context}
  {raoyehudayoff-ch01-deterministic-protocols}
  {pdf-fa052d8c95c0-p004-b005}
\uses{CommunicationComplexity.Deterministic.Protocol, CommunicationComplexity.Deterministic.Protocol.IsSubprotocol, CommunicationComplexity.Deterministic.Protocol.IsSubprotocol.trans, CommunicationComplexity.Deterministic.Protocol.numLeaves, CommunicationComplexity.Deterministic.Protocol.shape}
\difficulty{5}
Let $p$ be a deterministic two-party communication protocol, and let $n$ be a natural
number with $n > 1$. Write $\ell(q)$ for the number of leaves of a protocol $q$, and
suppose that $3\,\ell(p) \ge 2n$. Then there is a subprotocol $s$ of $p$ (a rooted
subtree of $p$) whose leaf count satisfies
\[
n \le 3\,\ell(s) < 2n.
\]
\end{lemma}

\begin{theorem}[Existence of a balanced subprotocol]
\lean{CommunicationComplexity.Deterministic.Protocol.balanced_subprotocol}
\proofsource{raoyehudayoff-ch01-deterministic-protocols}{
  pdf-fa052d8c95c0-p004-b003,
  pdf-fa052d8c95c0-p004-b004
}
\label{CommunicationComplexity.Deterministic.Protocol.balanced_subprotocol}
\leanok
\proofstep
  {CommunicationComplexity.Deterministic.Protocol}
  {direct}
  {raoyehudayoff-ch01-deterministic-protocols}
  {pdf-fa052d8c95c0-p003-b001}
\proofstep
  {CommunicationComplexity.Deterministic.Protocol.run}
  {direct}
  {raoyehudayoff-ch01-deterministic-protocols}
  {pdf-fa052d8c95c0-p003-b002}
\proofstep
  {CommunicationComplexity.Deterministic.Protocol.complexity}
  {direct}
  {raoyehudayoff-ch01-deterministic-protocols}
  {pdf-fa052d8c95c0-p003-b003}
\proofstep
  {CommunicationComplexity.Deterministic.Protocol.swap}
  {context}
  {raoyehudayoff-ch01-deterministic-protocols}
  {pdf-fa052d8c95c0-p003-b001}
\proofstep
  {CommunicationComplexity.Deterministic.Protocol.swap_run}
  {context}
  {raoyehudayoff-ch01-deterministic-protocols}
  {pdf-fa052d8c95c0-p003-b002}
\proofstep
  {CommunicationComplexity.Deterministic.Protocol.swap_complexity}
  {context}
  {raoyehudayoff-ch01-deterministic-protocols}
  {pdf-fa052d8c95c0-p003-b003}
\proofstep
  {CommunicationComplexity.Deterministic.Protocol.swap_swap}
  {context}
  {raoyehudayoff-ch01-deterministic-protocols}
  {pdf-fa052d8c95c0-p003-b001}
\proofstep
  {CommunicationComplexity.Deterministic.Protocol.comap}
  {context}
  {raoyehudayoff-ch01-deterministic-protocols}
  {pdf-fa052d8c95c0-p003-b001}
\proofstep
  {CommunicationComplexity.Deterministic.Protocol.comap_run}
  {context}
  {raoyehudayoff-ch01-deterministic-protocols}
  {pdf-fa052d8c95c0-p003-b002}
\proofstep
  {CommunicationComplexity.Deterministic.Protocol.comap_complexity}
  {context}
  {raoyehudayoff-ch01-deterministic-protocols}
  {pdf-fa052d8c95c0-p003-b003}
\proofstep
  {CommunicationComplexity.Deterministic.Protocol.shape}
  {context}
  {raoyehudayoff-ch01-deterministic-protocols}
  {pdf-fa052d8c95c0-p003-b005}
\proofstep
  {CommunicationComplexity.Deterministic.Protocol.numLeaves}
  {context}
  {raoyehudayoff-ch01-deterministic-protocols}
  {pdf-fa052d8c95c0-p003-b005,
    pdf-fa052d8c95c0-p004-b003}
\proofstep
  {CommunicationComplexity.Deterministic.Protocol.IsSubprotocol}
  {context}
  {raoyehudayoff-ch01-deterministic-protocols}
  {pdf-fa052d8c95c0-p004-b003}
\proofstep
  {CommunicationComplexity.Deterministic.Protocol.IsSubprotocol.trans}
  {context}
  {raoyehudayoff-ch01-deterministic-protocols}
  {pdf-fa052d8c95c0-p004-b004}
\proofstep
  {CommunicationComplexity.Deterministic.Protocol.balanced_aux}
  {direct}
  {raoyehudayoff-ch01-deterministic-protocols}
  {pdf-fa052d8c95c0-p004-b003,
    pdf-fa052d8c95c0-p004-b004}
\proofstep
  {CommunicationComplexity.Deterministic.Protocol.SubprotocolPath}
  {context}
  {raoyehudayoff-ch01-deterministic-protocols}
  {pdf-fa052d8c95c0-p004-b004}
\proofstep
  {CommunicationComplexity.Deterministic.Protocol.path_exists_of_isSubprotocol}
  {context}
  {raoyehudayoff-ch01-deterministic-protocols}
  {pdf-fa052d8c95c0-p004-b004}
\proofstep
  {CommunicationComplexity.Deterministic.Protocol.choosePath}
  {context}
  {raoyehudayoff-ch01-deterministic-protocols}
  {pdf-fa052d8c95c0-p004-b004}
\proofstep
  {CommunicationComplexity.Deterministic.Protocol.reachXPath}
  {context}
  {raoyehudayoff-ch01-deterministic-protocols}
  {pdf-fa052d8c95c0-p004-b005,
    pdf-fa052d8c95c0-p005-b002}
\proofstep
  {CommunicationComplexity.Deterministic.Protocol.reachYPath}
  {context}
  {raoyehudayoff-ch01-deterministic-protocols}
  {pdf-fa052d8c95c0-p004-b005,
    pdf-fa052d8c95c0-p005-b002}
\proofstep
  {CommunicationComplexity.Deterministic.Protocol.reachX}
  {context}
  {raoyehudayoff-ch01-deterministic-protocols}
  {pdf-fa052d8c95c0-p004-b005}
\proofstep
  {CommunicationComplexity.Deterministic.Protocol.reachY}
  {context}
  {raoyehudayoff-ch01-deterministic-protocols}
  {pdf-fa052d8c95c0-p004-b005}
\proofstep
  {CommunicationComplexity.Deterministic.Protocol.testSubprotocol}
  {context}
  {raoyehudayoff-ch01-deterministic-protocols}
  {pdf-fa052d8c95c0-p004-b005}
\proofstep
  {CommunicationComplexity.Deterministic.Protocol.testSubprotocol_complexity}
  {context}
  {raoyehudayoff-ch01-deterministic-protocols}
  {pdf-fa052d8c95c0-p004-b005}
\uses{CommunicationComplexity.Deterministic.Protocol, CommunicationComplexity.Deterministic.Protocol.IsSubprotocol, CommunicationComplexity.Deterministic.Protocol.IsSubprotocol.trans, CommunicationComplexity.Deterministic.Protocol.balanced_aux, CommunicationComplexity.Deterministic.Protocol.numLeaves, CommunicationComplexity.Deterministic.Protocol.shape}
\difficulty{5}
Let $p$ be a deterministic two-party communication protocol, and for any such protocol
$q$ write $\ell(q)$ for its number of leaves (its output nodes). Suppose $\ell(p) > 1$.
Then there is a protocol $s$ that is a rooted subtree of $p$ whose number of leaves
satisfies
\[
  \ell(p) \;\le\; 3\,\ell(s) \;<\; 2\,\ell(p),
\]
so that $s$ carries at least one-third and strictly fewer than two-thirds of the leaves
of $p$.
\end{theorem}

\begin{definition}[Subprotocol path witness]
\lean{CommunicationComplexity.Deterministic.Protocol.SubprotocolPath}
\label{CommunicationComplexity.Deterministic.Protocol.SubprotocolPath}
\leanok
\uses{CommunicationComplexity.Deterministic.Protocol}
The inductive type $\texttt{SubprotocolPath}\; s\; p$ is a data-carrying witness that $s$ is a
rooted subtree of $p$, recording the sequence of left/right choices made from the root of
$p$ down to the root of $s$.  It mirrors the constructors of \texttt{IsSubprotocol} but lives
in \texttt{Type} rather than \texttt{Prop}, enabling structural recursion over the path.
\end{definition}

\begin{definition}[Path to subprotocol proof]
\lean{CommunicationComplexity.Deterministic.Protocol.SubprotocolPath.toIsSubprotocol}
\label{CommunicationComplexity.Deterministic.Protocol.SubprotocolPath.toIsSubprotocol}
\leanok
\uses{CommunicationComplexity.Deterministic.Protocol, CommunicationComplexity.Deterministic.Protocol.IsSubprotocol, CommunicationComplexity.Deterministic.Protocol.SubprotocolPath}
Every \texttt{SubprotocolPath} $s\; p$ can be mapped to a proof that $s$ is a subprotocol of
$p$: the function \texttt{SubprotocolPath.toIsSubprotocol} performs this conversion by
structural recursion on the path.
\end{definition}

\begin{theorem}[Every subprotocol relation has a path witness]
\lean{CommunicationComplexity.Deterministic.Protocol.path_exists_of_isSubprotocol}
\label{CommunicationComplexity.Deterministic.Protocol.path_exists_of_isSubprotocol}
\leanok
\uses{CommunicationComplexity.Deterministic.Protocol, CommunicationComplexity.Deterministic.Protocol.IsSubprotocol, CommunicationComplexity.Deterministic.Protocol.SubprotocolPath, CommunicationComplexity.Deterministic.Protocol.reachXPath}
\difficulty{3}
Let $s$ and $p$ be deterministic two-party communication protocols over the same input
types $X$ and $Y$ and output type $\alpha$. If $s$ is a subprotocol of $p$ — that is,
the subprotocol predicate holds of the pair $(s, p)$ — then the type of subprotocol
paths from $s$ to $p$ is nonempty: there exists a data-carrying witness recording the
sequence of branch choices leading from the root of $p$ down to the root of $s$.
\end{theorem}

\begin{definition}[Classical choice of a subprotocol path]
\lean{CommunicationComplexity.Deterministic.Protocol.choosePath}
\label{CommunicationComplexity.Deterministic.Protocol.choosePath}
\leanok
\uses{CommunicationComplexity.Deterministic.Protocol, CommunicationComplexity.Deterministic.Protocol.IsSubprotocol, CommunicationComplexity.Deterministic.Protocol.SubprotocolPath, CommunicationComplexity.Deterministic.Protocol.path_exists_of_isSubprotocol}
Given a proof that $s$ is a subprotocol of $p$, \texttt{choosePath} uses classical choice to
extract a definite \texttt{SubprotocolPath} $s\; p$ witness.  The resulting function is
noncomputable.
\end{definition}

\begin{lemma}[Leaf count is monotone along subprotocol paths]
\lean{CommunicationComplexity.Deterministic.Protocol.SubprotocolPath.numLeaves_le}
\label{CommunicationComplexity.Deterministic.Protocol.SubprotocolPath.numLeaves_le}
\leanok
\uses{CommunicationComplexity.Deterministic.Protocol, CommunicationComplexity.Deterministic.Protocol.SubprotocolPath, CommunicationComplexity.Deterministic.Protocol.numLeaves, CommunicationComplexity.Deterministic.Protocol.shape}
\difficulty{4}
Let $s$ and $p$ be deterministic two-party communication protocols over the same input
types, and suppose there is a subprotocol path from $s$ to $p$, witnessing that $s$
occurs as a rooted subtree of $p$. Then $s$ has at most as many leaves as $p$; that is,
the number of leaves of $s$ is less than or equal to the number of leaves of $p$.
\end{lemma}

\begin{definition}[Reachable Alice inputs along a path]
\lean{CommunicationComplexity.Deterministic.Protocol.reachXPath}
\statementsource{raoyehudayoff-ch01-deterministic-protocols}{pdf-fa052d8c95c0-p005-b002}
\label{CommunicationComplexity.Deterministic.Protocol.reachXPath}
\leanok
\uses{CommunicationComplexity.Deterministic.Protocol, CommunicationComplexity.Deterministic.Protocol.SubprotocolPath}
Given a subprotocol path \texttt{hsp} from $s$ to $p$, \texttt{reachXPath hsp} is the set of
Alice's inputs $x \in X$ that are consistent with all of Alice's choices along the path from
the root of $p$ to the root of $s$.
\end{definition}

\begin{definition}[Reachable Bob inputs along a path]
\lean{CommunicationComplexity.Deterministic.Protocol.reachYPath}
\statementsource{raoyehudayoff-ch01-deterministic-protocols}{pdf-fa052d8c95c0-p005-b002}
\label{CommunicationComplexity.Deterministic.Protocol.reachYPath}
\leanok
\uses{CommunicationComplexity.Deterministic.Protocol, CommunicationComplexity.Deterministic.Protocol.SubprotocolPath}
Given a subprotocol path \texttt{hsp} from $s$ to $p$, \texttt{reachYPath hsp} is the set of
Bob's inputs $y \in Y$ that are consistent with all of Bob's choices along the path from the
root of $p$ to the root of $s$.
\end{definition}

\begin{definition}[Inputs that reach a subprotocol path]
\lean{CommunicationComplexity.Deterministic.Protocol.reachesPath}
\statementsource{raoyehudayoff-ch01-deterministic-protocols}{pdf-fa052d8c95c0-p005-b002}
\label{CommunicationComplexity.Deterministic.Protocol.reachesPath}
\leanok
\uses{CommunicationComplexity.Deterministic.Protocol, CommunicationComplexity.Deterministic.Protocol.SubprotocolPath, CommunicationComplexity.Deterministic.Protocol.reachXPath, CommunicationComplexity.Deterministic.Protocol.reachYPath}
$\texttt{reachesPath}\;\texttt{hsp}\;x\;y$ is the proposition that the input pair $(x, y)$
causes protocol $p$ to reach the subprotocol $s$ indicated by the path \texttt{hsp}; that is,
$x \in \texttt{reachXPath}\;\texttt{hsp}$ and $y \in \texttt{reachYPath}\;\texttt{hsp}$.
\end{definition}

\begin{definition}[Alice's reach set for a subprotocol]
\lean{CommunicationComplexity.Deterministic.Protocol.reachX}
\statementsource{raoyehudayoff-ch01-deterministic-protocols}{pdf-fa052d8c95c0-p005-b002}
\label{CommunicationComplexity.Deterministic.Protocol.reachX}
\leanok
\uses{CommunicationComplexity.Deterministic.Protocol, CommunicationComplexity.Deterministic.Protocol.IsSubprotocol, CommunicationComplexity.Deterministic.Protocol.choosePath, CommunicationComplexity.Deterministic.Protocol.reachXPath}
Given a proof \texttt{hsp} that $s$ is a subprotocol of $p$, \texttt{reachX hsp} is the set of
Alice's inputs that reach the subprotocol, defined by applying \texttt{reachXPath} to the
classically chosen path.
\end{definition}

\begin{definition}[Bob's reach set for a subprotocol]
\lean{CommunicationComplexity.Deterministic.Protocol.reachY}
\statementsource{raoyehudayoff-ch01-deterministic-protocols}{pdf-fa052d8c95c0-p005-b002}
\label{CommunicationComplexity.Deterministic.Protocol.reachY}
\leanok
\uses{CommunicationComplexity.Deterministic.Protocol, CommunicationComplexity.Deterministic.Protocol.IsSubprotocol, CommunicationComplexity.Deterministic.Protocol.choosePath, CommunicationComplexity.Deterministic.Protocol.reachYPath}
Given a proof \texttt{hsp} that $s$ is a subprotocol of $p$, \texttt{reachY hsp} is the set of
Bob's inputs that reach the subprotocol, defined by applying \texttt{reachYPath} to the
classically chosen path.
\end{definition}

\begin{definition}[Inputs that reach a subprotocol]
\lean{CommunicationComplexity.Deterministic.Protocol.reaches}
\statementsource{raoyehudayoff-ch01-deterministic-protocols}{pdf-fa052d8c95c0-p005-b002}
\label{CommunicationComplexity.Deterministic.Protocol.reaches}
\leanok
\uses{CommunicationComplexity.Deterministic.Protocol, CommunicationComplexity.Deterministic.Protocol.IsSubprotocol, CommunicationComplexity.Deterministic.Protocol.reachX, CommunicationComplexity.Deterministic.Protocol.reachY}
$\texttt{reaches}\;\texttt{hsp}\;x\;y$ is the proposition that input $(x, y)$ causes $p$ to
reach the subprotocol $s$; equivalently, $x \in \texttt{reachX}\;\texttt{hsp}$ and
$y \in \texttt{reachY}\;\texttt{hsp}$.
\end{definition}

\begin{lemma}[Running a protocol and a reached subprotocol agree]
\lean{CommunicationComplexity.Deterministic.Protocol.subprotocol_run_eq_of_reachesPath}
\label{CommunicationComplexity.Deterministic.Protocol.subprotocol_run_eq_of_reachesPath}
\leanok
\uses{CommunicationComplexity.Deterministic.Protocol, CommunicationComplexity.Deterministic.Protocol.SubprotocolPath, CommunicationComplexity.Deterministic.Protocol.reachXPath, CommunicationComplexity.Deterministic.Protocol.reachYPath, CommunicationComplexity.Deterministic.Protocol.reachesPath, CommunicationComplexity.Deterministic.Protocol.run}
\difficulty{4}
Let $p$ and $s$ be deterministic two-party communication protocols over Alice's input
type $X$ and Bob's input type $Y$ with output type $\alpha$, and let $\pi$ be a
subprotocol path witnessing that $s$ sits at a rooted subtree of $p$. Suppose the input
pair $(x, y)$ reaches $s$ along $\pi$, meaning $x$ lies in the set of Alice inputs
consistent with all of Alice's branching choices on $\pi$ and $y$ lies in the set of Bob
inputs consistent with all of Bob's branching choices on $\pi$. Then executing $p$ and
executing $s$ on $(x, y)$ yield the same output.
\end{lemma}

\begin{lemma}[Subprotocol and parent agree on inputs reaching it]
\lean{CommunicationComplexity.Deterministic.Protocol.subprotocol_run_eq_of_reaches}
\label{CommunicationComplexity.Deterministic.Protocol.subprotocol_run_eq_of_reaches}
\leanok
\uses{CommunicationComplexity.Deterministic.Protocol, CommunicationComplexity.Deterministic.Protocol.IsSubprotocol, CommunicationComplexity.Deterministic.Protocol.choosePath, CommunicationComplexity.Deterministic.Protocol.reachX, CommunicationComplexity.Deterministic.Protocol.reachY, CommunicationComplexity.Deterministic.Protocol.reaches, CommunicationComplexity.Deterministic.Protocol.run, CommunicationComplexity.Deterministic.Protocol.subprotocol_run_eq_of_reachesPath}
\difficulty{1}
Let $p$ and $s$ be deterministic two-party communication protocols with Alice's inputs
drawn from $X$, Bob's inputs from $Y$, and outputs in $\alpha$, and suppose $s$ is a
subprotocol of $p$, that is, a rooted subtree of $p$. If an input pair $(x, y)$ reaches
$s$ — meaning $x$ is consistent with every one of Alice's branch choices and $y$ with
every one of Bob's branch choices along the path from the root of $p$ down to the root
of $s$ — then evaluating $p$ on $(x, y)$ returns the same output as evaluating $s$ on
$(x, y)$.
\end{lemma}

\begin{definition}[Canonical output of a protocol]
\lean{CommunicationComplexity.Deterministic.Protocol.chooseOutput}
\label{CommunicationComplexity.Deterministic.Protocol.chooseOutput}
\leanok
\uses{CommunicationComplexity.Deterministic.Protocol}
\texttt{chooseOutput p} returns a canonical output value from protocol $p$ by always following
the false (left) branch at every Alice-node; it yields the value at the leftmost leaf.
\end{definition}

\begin{definition}[Erase a subprotocol path]
\lean{CommunicationComplexity.Deterministic.Protocol.erasePath}
\label{CommunicationComplexity.Deterministic.Protocol.erasePath}
\leanok
\uses{CommunicationComplexity.Deterministic.Protocol, CommunicationComplexity.Deterministic.Protocol.SubprotocolPath, CommunicationComplexity.Deterministic.Protocol.chooseOutput}
Given a subprotocol path \texttt{hsp} from $s$ to $p$, \texttt{erasePath hsp} produces a new
protocol that is $p$ with the subtree $s$ collapsed to a single leaf whose value is
$\texttt{chooseOutput}\;s$ (the leftmost leaf of $s$).
\end{definition}

\begin{lemma}[Leaf count after collapsing a subprotocol]
\lean{CommunicationComplexity.Deterministic.Protocol.erasePath_numLeaves}
\label{CommunicationComplexity.Deterministic.Protocol.erasePath_numLeaves}
\leanok
\uses{CommunicationComplexity.Deterministic.Protocol, CommunicationComplexity.Deterministic.Protocol.SubprotocolPath, CommunicationComplexity.Deterministic.Protocol.SubprotocolPath.numLeaves_le, CommunicationComplexity.Deterministic.Protocol.erasePath, CommunicationComplexity.Deterministic.Protocol.numLeaves, CommunicationComplexity.Deterministic.Protocol.shape}
\difficulty{5}
Let $p$ be a deterministic two-party communication protocol, and let $s$ be a
subprotocol of $p$, witnessed by a subprotocol path recording the sequence of branch
choices leading from the root of $p$ down to the root of $s$. Form the protocol obtained
from $p$ by collapsing the entire subtree $s$ to a single output leaf carrying its
canonical output value. Writing $\ell(q)$ for the number of leaves (output nodes) of a
protocol $q$, the collapsed protocol has $\ell(p) - \ell(s) + 1$ leaves, where the
subtraction is truncated subtraction of natural numbers.
\end{lemma}

\begin{lemma}[Erased subprotocol agrees off the reachable region]
\lean{CommunicationComplexity.Deterministic.Protocol.erasePath_run_outside}
\label{CommunicationComplexity.Deterministic.Protocol.erasePath_run_outside}
\leanok
\uses{CommunicationComplexity.Deterministic.Protocol, CommunicationComplexity.Deterministic.Protocol.SubprotocolPath, CommunicationComplexity.Deterministic.Protocol.erasePath, CommunicationComplexity.Deterministic.Protocol.reachXPath, CommunicationComplexity.Deterministic.Protocol.reachYPath, CommunicationComplexity.Deterministic.Protocol.reachesPath, CommunicationComplexity.Deterministic.Protocol.run}
\difficulty{5}
Let $p$ be a deterministic two-party communication protocol over Alice's inputs $X$ and
Bob's inputs $Y$, let $s$ be a rooted subprotocol of $p$ recorded by a subprotocol path
from $s$ to $p$, and let $p'$ be the protocol obtained from $p$ by collapsing the
subtree $s$ to a single leaf carrying the canonical output of $s$ (its leftmost leaf
value). Suppose an input pair $(x,y)$, with $x \in X$ and $y \in Y$, does not reach $s$:
it fails to be consistent with all of Alice's and Bob's branch choices along the path
from the root of $p$ down to the root of $s$. Then running $p'$ on $(x,y)$ produces the
same output as running $p$ on $(x,y)$.
\end{lemma}

\begin{definition}[Erase a subprotocol (propositional version)]
\lean{CommunicationComplexity.Deterministic.Protocol.erase}
\label{CommunicationComplexity.Deterministic.Protocol.erase}
\leanok
\uses{CommunicationComplexity.Deterministic.Protocol, CommunicationComplexity.Deterministic.Protocol.IsSubprotocol, CommunicationComplexity.Deterministic.Protocol.choosePath, CommunicationComplexity.Deterministic.Protocol.erasePath}
The noncomputable variant of \texttt{erasePath} that takes an \texttt{IsSubprotocol} proof
instead of a path: \texttt{erase hsp} applies \texttt{erasePath} to the classically chosen path
for \texttt{hsp}.
\end{definition}

\begin{lemma}[Leaf count after erasing a subprotocol]
\lean{CommunicationComplexity.Deterministic.Protocol.erase_numLeaves}
\label{CommunicationComplexity.Deterministic.Protocol.erase_numLeaves}
\leanok
\uses{CommunicationComplexity.Deterministic.Protocol, CommunicationComplexity.Deterministic.Protocol.IsSubprotocol, CommunicationComplexity.Deterministic.Protocol.choosePath, CommunicationComplexity.Deterministic.Protocol.erase, CommunicationComplexity.Deterministic.Protocol.erasePath_numLeaves, CommunicationComplexity.Deterministic.Protocol.numLeaves}
\difficulty{1}
Let $s$ and $p$ be deterministic two-party communication protocols over the same input
types $X$, $Y$ and output type $\alpha$, and suppose that $s$ is a subprotocol of $p$,
i.e.\ a rooted subtree of $p$. Write $L(q)$ for the number of leaves (output nodes) of a
protocol $q$. Erasing $s$ from $p$ — collapsing the subtree $s$ to a single leaf
carrying its canonical (leftmost) output value — yields a protocol whose number of
leaves is
\[
L(p) - L(s) + 1,
\]
where the subtraction is truncated natural-number subtraction.
\end{lemma}

\begin{lemma}[Erasing a subprotocol preserves behaviour off its reach set]
\lean{CommunicationComplexity.Deterministic.Protocol.erase_run_outside}
\label{CommunicationComplexity.Deterministic.Protocol.erase_run_outside}
\leanok
\uses{CommunicationComplexity.Deterministic.Protocol, CommunicationComplexity.Deterministic.Protocol.IsSubprotocol, CommunicationComplexity.Deterministic.Protocol.choosePath, CommunicationComplexity.Deterministic.Protocol.erase, CommunicationComplexity.Deterministic.Protocol.erasePath_run_outside, CommunicationComplexity.Deterministic.Protocol.reachX, CommunicationComplexity.Deterministic.Protocol.reachY, CommunicationComplexity.Deterministic.Protocol.reaches, CommunicationComplexity.Deterministic.Protocol.run}
\difficulty{1}
Let $p$ be a deterministic two-party communication protocol, and let $s$ be a
subprotocol of $p$ (a rooted subtree of $p$). Form the erased protocol
$\mathrm{erase}(s,p)$, obtained from $p$ by collapsing the subtree $s$ to a single leaf.
Then for every pair of inputs $x$ (Alice) and $y$ (Bob) that does not reach $s$, running
the erased protocol and running $p$ produce the same output:
\[
\mathrm{run}\bigl(\mathrm{erase}(s,p)\bigr)(x,y) = \mathrm{run}(p)(x,y).
\]
\end{lemma}

\begin{definition}[Delete a subprotocol path by splicing]
\lean{CommunicationComplexity.Deterministic.Protocol.deletePath}
\statementsource{raoyehudayoff-ch01-deterministic-protocols}{pdf-fa052d8c95c0-p004-b005}
\label{CommunicationComplexity.Deterministic.Protocol.deletePath}
\leanok
\uses{CommunicationComplexity.Deterministic.Protocol, CommunicationComplexity.Deterministic.Protocol.SubprotocolPath}
\texttt{deletePath hsp} removes the subtree $s$ from $p$ by splicing out the branch leading to
$s$; the result is \texttt{none} if $s = p$ (the root case) and \texttt{some q} otherwise, where
$q$ is $p$ with the branch to $s$ deleted.
\end{definition}

\begin{lemma}[Deleting a strict subprotocol yields a protocol]
\lean{CommunicationComplexity.Deterministic.Protocol.deletePath_ne_none_of_lt}
\label{CommunicationComplexity.Deterministic.Protocol.deletePath_ne_none_of_lt}
\leanok
\uses{CommunicationComplexity.Deterministic.Protocol, CommunicationComplexity.Deterministic.Protocol.SubprotocolPath, CommunicationComplexity.Deterministic.Protocol.deletePath, CommunicationComplexity.Deterministic.Protocol.numLeaves}
\difficulty{3}
Let $s$ and $p$ be deterministic two-party communication protocols, and suppose that $s$
occurs as a rooted subprotocol of $p$, witnessed by a path recording the sequence of
left/right branchings from the root of $p$ down to the root of $s$. If $s$ has strictly
fewer leaves (output nodes) than $p$, then splicing out the branch leading to $s$ does
not collapse the whole protocol: the result of the deletion is some protocol rather than
the empty result.
\end{lemma}

\begin{lemma}[Deletion succeeds along a proper subprotocol path]
\lean{CommunicationComplexity.Deterministic.Protocol.deletePath_exists_of_lt}
\label{CommunicationComplexity.Deterministic.Protocol.deletePath_exists_of_lt}
\leanok
\uses{CommunicationComplexity.Deterministic.Protocol, CommunicationComplexity.Deterministic.Protocol.SubprotocolPath, CommunicationComplexity.Deterministic.Protocol.deletePath, CommunicationComplexity.Deterministic.Protocol.deletePath_ne_none_of_lt, CommunicationComplexity.Deterministic.Protocol.numLeaves}
\difficulty{2}
Let $s$ and $p$ be deterministic two-party communication protocols, and suppose $s$
occurs as a rooted subtree of $p$, witnessed by a subprotocol path recording the
branching choices from the root of $p$ down to the root of $s$. If $s$ has strictly
fewer leaves than $p$, then deleting the branch leading to $s$ yields a protocol: the
splicing operation along this path is defined and produces some protocol $q$.
\end{lemma}

\begin{definition}[Prune a subprotocol path]
\lean{CommunicationComplexity.Deterministic.Protocol.prunePath}
\label{CommunicationComplexity.Deterministic.Protocol.prunePath}
\leanok
\uses{CommunicationComplexity.Deterministic.Protocol, CommunicationComplexity.Deterministic.Protocol.SubprotocolPath, CommunicationComplexity.Deterministic.Protocol.deletePath_exists_of_lt, CommunicationComplexity.Deterministic.Protocol.numLeaves}
Given a path \texttt{hsp} from $s$ to $p$ with $s.\mathrm{numLeaves} < p.\mathrm{numLeaves}$,
\texttt{prunePath hsp hlt} is the protocol obtained by classically extracting the \texttt{some}
value from \texttt{deletePath hsp}.
\end{definition}

\begin{lemma}[Path deletion yields the pruned protocol]
\lean{CommunicationComplexity.Deterministic.Protocol.prunePath_spec}
\label{CommunicationComplexity.Deterministic.Protocol.prunePath_spec}
\leanok
\uses{CommunicationComplexity.Deterministic.Protocol, CommunicationComplexity.Deterministic.Protocol.SubprotocolPath, CommunicationComplexity.Deterministic.Protocol.deletePath, CommunicationComplexity.Deterministic.Protocol.deletePath_exists_of_lt, CommunicationComplexity.Deterministic.Protocol.numLeaves, CommunicationComplexity.Deterministic.Protocol.prunePath}
\difficulty{1}
Let $p$ be a deterministic communication protocol and let $s$ be a rooted subprotocol of
$p$, recorded by a subprotocol path running from the root of $p$ down to the root of
$s$. Suppose that $s$ has strictly fewer leaves than $p$. Then splicing this path out of
$p$ does not fall into the empty root case, in which the subprotocol coincides with the
whole protocol, but instead returns a genuine protocol; and that protocol is exactly the
pruned protocol determined by the path together with this leaf-count hypothesis.
\end{lemma}

\begin{lemma}[Deleting a subprotocol returns nothing only for the whole protocol]
\lean{CommunicationComplexity.Deterministic.Protocol.eq_of_deletePath_none}
\label{CommunicationComplexity.Deterministic.Protocol.eq_of_deletePath_none}
\leanok
\uses{CommunicationComplexity.Deterministic.Protocol, CommunicationComplexity.Deterministic.Protocol.SubprotocolPath, CommunicationComplexity.Deterministic.Protocol.deletePath}
\difficulty{3}
Let $s$ and $p$ be deterministic two-party communication protocols over input types $X$
(Alice) and $Y$ (Bob) producing values in $\alpha$, and suppose $s$ is a rooted subtree
of $p$, as recorded by a subprotocol path from the root of $p$ down to the root of $s$.
If splicing that path out of $p$ produces no protocol at all, then $s$ is in fact all of
$p$; that is, $s = p$.
\end{lemma}

\begin{lemma}[Leaf count after deleting a subprotocol branch]
\lean{CommunicationComplexity.Deterministic.Protocol.deletePath_numLeaves_of_some}
\label{CommunicationComplexity.Deterministic.Protocol.deletePath_numLeaves_of_some}
\leanok
\uses{CommunicationComplexity.Deterministic.Protocol, CommunicationComplexity.Deterministic.Protocol.SubprotocolPath, CommunicationComplexity.Deterministic.Protocol.SubprotocolPath.numLeaves_le, CommunicationComplexity.Deterministic.Protocol.deletePath, CommunicationComplexity.Deterministic.Protocol.eq_of_deletePath_none, CommunicationComplexity.Deterministic.Protocol.numLeaves, CommunicationComplexity.Deterministic.Protocol.shape}
\difficulty{5}
Let $p$ be a deterministic communication protocol and let $s$ be a rooted subtree of
$p$, witnessed by a subprotocol path from the root of $p$ down to the root of $s$. Write
$\ell(r)$ for the number of leaves (output nodes) of a protocol $r$. If deleting this
branch splices out a proper subtree — that is, if the deletion produces a protocol $q$
rather than the empty result of the root case — then \[ \ell(q) = \ell(p) - \ell(s), \]
where the subtraction is truncated subtraction on $\bbn$.
\end{lemma}

\begin{lemma}[Leaf count after pruning a subprotocol path]
\lean{CommunicationComplexity.Deterministic.Protocol.prunePath_numLeaves_of_lt}
\label{CommunicationComplexity.Deterministic.Protocol.prunePath_numLeaves_of_lt}
\leanok
\uses{CommunicationComplexity.Deterministic.Protocol, CommunicationComplexity.Deterministic.Protocol.SubprotocolPath, CommunicationComplexity.Deterministic.Protocol.deletePath_numLeaves_of_some, CommunicationComplexity.Deterministic.Protocol.numLeaves, CommunicationComplexity.Deterministic.Protocol.prunePath, CommunicationComplexity.Deterministic.Protocol.prunePath_spec}
\difficulty{1}
Let $s$ and $p$ be deterministic two-party communication protocols, and let $s$ lie
along a subprotocol path inside $p$, so that $s$ occurs as a rooted subtree of $p$. If
$s$ has strictly fewer leaves than $p$, that is, if the number of leaves of $s$ is less
than the number of leaves of $p$, then the protocol obtained by pruning this path from
$p$ has exactly $\text{numLeaves}(p) - \text{numLeaves}(s)$ leaves.
\end{lemma}

\begin{lemma}[Trivial path is reached by every input]
\lean{CommunicationComplexity.Deterministic.Protocol.reachesPath_of_deletePath_none}
\label{CommunicationComplexity.Deterministic.Protocol.reachesPath_of_deletePath_none}
\leanok
\uses{CommunicationComplexity.Deterministic.Protocol, CommunicationComplexity.Deterministic.Protocol.SubprotocolPath, CommunicationComplexity.Deterministic.Protocol.deletePath, CommunicationComplexity.Deterministic.Protocol.reachXPath, CommunicationComplexity.Deterministic.Protocol.reachYPath, CommunicationComplexity.Deterministic.Protocol.reachesPath}
\difficulty{3}
Let $p$ be a deterministic two-party communication protocol with Alice inputs in $X$ and
Bob inputs in $Y$, and let $\pi$ be a subprotocol path witnessing that a subprotocol $s$
is a rooted subtree of $p$. If deleting the path $\pi$ from $p$ returns no protocol —
which occurs precisely when the path is trivial, so that $s = p$ — then every input pair
$(x, y) \in X \times Y$ reaches $s$ along $\pi$, meaning $x$ is consistent with all of
Alice's choices along $\pi$ and $y$ is consistent with all of Bob's choices along $\pi$.
\end{lemma}

\begin{lemma}[Deleting an unreached subprotocol preserves outputs]
\lean{CommunicationComplexity.Deterministic.Protocol.deletePath_run_outside_of_some}
\label{CommunicationComplexity.Deterministic.Protocol.deletePath_run_outside_of_some}
\leanok
\uses{CommunicationComplexity.Deterministic.Protocol, CommunicationComplexity.Deterministic.Protocol.SubprotocolPath, CommunicationComplexity.Deterministic.Protocol.deletePath, CommunicationComplexity.Deterministic.Protocol.reachXPath, CommunicationComplexity.Deterministic.Protocol.reachYPath, CommunicationComplexity.Deterministic.Protocol.reachesPath, CommunicationComplexity.Deterministic.Protocol.reachesPath_of_deletePath_none, CommunicationComplexity.Deterministic.Protocol.run}
\difficulty{5}
Let $p$ be a deterministic two-party communication protocol on input types $X$ (Alice)
and $Y$ (Bob), and let $s$ be a rooted subprotocol of $p$ witnessed by a subprotocol
path from $s$ to $p$. Suppose that deleting this path — splicing out the branch of $p$
leading to $s$ — yields a protocol $q$ (rather than being undefined, which happens only
when $s$ is the whole of $p$). Then for every pair of inputs $(x,y)$ that does not reach
the subprotocol $s$ along the path — that is, $x$ is not a reachable Alice input or $y$
is not a reachable Bob input for that path — executing $q$ on $(x,y)$ produces the same
output as executing $p$ on $(x,y)$.
\end{lemma}

\begin{lemma}[Pruned protocol agrees off the pruned path]
\lean{CommunicationComplexity.Deterministic.Protocol.prunePath_run_outside_of_lt}
\label{CommunicationComplexity.Deterministic.Protocol.prunePath_run_outside_of_lt}
\leanok
\uses{CommunicationComplexity.Deterministic.Protocol, CommunicationComplexity.Deterministic.Protocol.SubprotocolPath, CommunicationComplexity.Deterministic.Protocol.deletePath_run_outside_of_some, CommunicationComplexity.Deterministic.Protocol.numLeaves, CommunicationComplexity.Deterministic.Protocol.prunePath, CommunicationComplexity.Deterministic.Protocol.prunePath_spec, CommunicationComplexity.Deterministic.Protocol.reachesPath, CommunicationComplexity.Deterministic.Protocol.run}
\difficulty{1}
Let $p$ be a deterministic two-party communication protocol with Alice inputs in $X$ and
Bob inputs in $Y$, and let $s$ be a rooted subprotocol of $p$, witnessed by a
subprotocol path from the root of $p$ down to the root of $s$. Suppose $s$ has strictly
fewer leaves than $p$, so that pruning this path yields a protocol $p'$. Then for every
input pair $(x,y)$ that does not reach $s$ along this path — that is, $x$ is not among
the Alice inputs consistent with all of Alice's branching choices on the path, or $y$ is
not among the Bob inputs consistent with all of Bob's choices on it — the pruned
protocol $p'$ and the original protocol $p$ return the same value on $(x,y)$.
\end{lemma}

\begin{definition}[Prune a subprotocol (propositional version)]
\lean{CommunicationComplexity.Deterministic.Protocol.prune}
\label{CommunicationComplexity.Deterministic.Protocol.prune}
\leanok
\uses{CommunicationComplexity.Deterministic.Protocol, CommunicationComplexity.Deterministic.Protocol.IsSubprotocol, CommunicationComplexity.Deterministic.Protocol.choosePath, CommunicationComplexity.Deterministic.Protocol.numLeaves, CommunicationComplexity.Deterministic.Protocol.prunePath}
The noncomputable variant of \texttt{prunePath} taking an \texttt{IsSubprotocol} proof:
\texttt{prune hsp hlt} applies \texttt{CommunicationComplexity.Deterministic.Protocol.prunePath} to the classically chosen path for
\texttt{hsp}, given the strictness condition $s.\mathrm{numLeaves} < p.\mathrm{numLeaves}$.
\end{definition}

\begin{lemma}[Leaf count after pruning a subprotocol]
\lean{CommunicationComplexity.Deterministic.Protocol.prune_numLeaves_of_lt}
\label{CommunicationComplexity.Deterministic.Protocol.prune_numLeaves_of_lt}
\leanok
\uses{CommunicationComplexity.Deterministic.Protocol, CommunicationComplexity.Deterministic.Protocol.IsSubprotocol, CommunicationComplexity.Deterministic.Protocol.choosePath, CommunicationComplexity.Deterministic.Protocol.numLeaves, CommunicationComplexity.Deterministic.Protocol.prune, CommunicationComplexity.Deterministic.Protocol.prunePath_numLeaves_of_lt}
\difficulty{1}
Let $p$ be a deterministic two-party communication protocol and let $s$ be a subprotocol
of $p$, that is, a rooted subtree of $p$. Write $L(q)$ for the number of leaves (output
nodes) of a protocol $q$, and suppose that $L(s) < L(p)$. Then the protocol obtained by
pruning $s$ out of $p$ has exactly $L(p) - L(s)$ leaves.
\end{lemma}

\begin{lemma}[Pruning preserves outputs on inputs that miss the subprotocol]
\lean{CommunicationComplexity.Deterministic.Protocol.prune_run_outside_of_lt}
\label{CommunicationComplexity.Deterministic.Protocol.prune_run_outside_of_lt}
\leanok
\uses{CommunicationComplexity.Deterministic.Protocol, CommunicationComplexity.Deterministic.Protocol.IsSubprotocol, CommunicationComplexity.Deterministic.Protocol.choosePath, CommunicationComplexity.Deterministic.Protocol.numLeaves, CommunicationComplexity.Deterministic.Protocol.prune, CommunicationComplexity.Deterministic.Protocol.prunePath_run_outside_of_lt, CommunicationComplexity.Deterministic.Protocol.reachX, CommunicationComplexity.Deterministic.Protocol.reachY, CommunicationComplexity.Deterministic.Protocol.reaches, CommunicationComplexity.Deterministic.Protocol.run}
\difficulty{1}
Let $s$ and $p$ be deterministic two-party communication protocols over Alice's inputs
$X$, Bob's inputs $Y$, and output type $\alpha$, and suppose $s$ is a subprotocol of $p$
(a rooted subtree) whose leaf count is strictly smaller than that of $p$. Form the
pruned protocol obtained by deleting the subtree $s$ from $p$. Then for every pair of
inputs $x \in X$ and $y \in Y$ that does not reach $s$ — that is, $(x,y)$ fails to be
consistent with all of the branching choices along the path from the root of $p$ down to
$s$ — the pruned protocol and $p$ produce the same output on $(x,y)$: running the pruned
protocol on $x$ and $y$ yields exactly $p$'s output on $x$ and $y$.
\end{lemma}

\begin{definition}[Test subprotocol construction]
\lean{CommunicationComplexity.Deterministic.Protocol.testSubprotocol}
\statementsource{raoyehudayoff-ch01-deterministic-protocols}{pdf-fa052d8c95c0-p004-b005}
\label{CommunicationComplexity.Deterministic.Protocol.testSubprotocol}
\leanok
\uses{CommunicationComplexity.Deterministic.Protocol, CommunicationComplexity.Deterministic.Protocol.IsSubprotocol, CommunicationComplexity.Deterministic.Protocol.reachX, CommunicationComplexity.Deterministic.Protocol.reachY}
Given that $s$ is a subprotocol of $p$ and two protocols \texttt{qIn} and \texttt{qOut},
\texttt{testSubprotocol hsp qIn qOut} is a noncomputable protocol that uses two bits of
communication to determine whether the input reaches $s$, then routes to \texttt{qIn} if it
does and to \texttt{qOut} otherwise.
\end{definition}

\begin{lemma}[Complexity of the test subprotocol]
\lean{CommunicationComplexity.Deterministic.Protocol.testSubprotocol_complexity}
\label{CommunicationComplexity.Deterministic.Protocol.testSubprotocol_complexity}
\leanok
\uses{CommunicationComplexity.Deterministic.Protocol, CommunicationComplexity.Deterministic.Protocol.IsSubprotocol, CommunicationComplexity.Deterministic.Protocol.complexity, CommunicationComplexity.Deterministic.Protocol.testSubprotocol}
\difficulty{2}
Let $X$ and $Y$ be the input types of Alice and Bob in a deterministic two-party
communication protocol producing values of type $\alpha$, and suppose $s$ is a
subprotocol of such a protocol $p$. For any two further protocols $q_{\mathrm{in}}$ and
$q_{\mathrm{out}}$, the test subprotocol — which spends one bit for Alice and one bit
for Bob to decide whether the input reaches $s$, then continues as $q_{\mathrm{in}}$ if
it does and as $q_{\mathrm{out}}$ otherwise — has communication complexity \[ 2 +
\max\bigl(\mathrm{complexity}(q_{\mathrm{in}}),\,
\mathrm{complexity}(q_{\mathrm{out}})\bigr). \]
\end{lemma}

\begin{lemma}[Test subprotocol runs the inside branch on reaching inputs]
\lean{CommunicationComplexity.Deterministic.Protocol.testSubprotocol_run_inside}
\label{CommunicationComplexity.Deterministic.Protocol.testSubprotocol_run_inside}
\leanok
\uses{CommunicationComplexity.Deterministic.Protocol, CommunicationComplexity.Deterministic.Protocol.IsSubprotocol, CommunicationComplexity.Deterministic.Protocol.reaches, CommunicationComplexity.Deterministic.Protocol.run, CommunicationComplexity.Deterministic.Protocol.testSubprotocol}
\difficulty{2}
Let $s$ be a subprotocol of a deterministic communication protocol $p$, and let
$q_{\mathrm{in}}$ and $q_{\mathrm{out}}$ be two further protocols on the same input
types. If the input pair $(x,y)$ reaches $s$ — that is, $x$ lies in Alice's reach set
and $y$ in Bob's reach set for $s$ — then executing the test subprotocol built from $s$,
$q_{\mathrm{in}}$, and $q_{\mathrm{out}}$ on $(x,y)$ yields the same output as executing
$q_{\mathrm{in}}$ on $(x,y)$.
\end{lemma}

\begin{lemma}[Test subprotocol agrees with the reject branch off the reach set]
\lean{CommunicationComplexity.Deterministic.Protocol.testSubprotocol_run_outside}
\label{CommunicationComplexity.Deterministic.Protocol.testSubprotocol_run_outside}
\leanok
\uses{CommunicationComplexity.Deterministic.Protocol, CommunicationComplexity.Deterministic.Protocol.IsSubprotocol, CommunicationComplexity.Deterministic.Protocol.reachX, CommunicationComplexity.Deterministic.Protocol.reachY, CommunicationComplexity.Deterministic.Protocol.reaches, CommunicationComplexity.Deterministic.Protocol.run, CommunicationComplexity.Deterministic.Protocol.testSubprotocol}
\difficulty{3}
Let $p$ be a deterministic two-party communication protocol over Alice's input type $X$
and Bob's input type $Y$ with outputs in $\alpha$, let $s$ be a subprotocol of $p$, and
let $q_{\mathrm{in}}$ and $q_{\mathrm{out}}$ be two further such protocols. Suppose an
input pair $(x, y) \in X \times Y$ does not reach $s$; that is, $x$ fails to lie in
Alice's reach set for $s$, or $y$ fails to lie in Bob's reach set for $s$. Then
executing the test subprotocol built from $s$, $q_{\mathrm{in}}$, and $q_{\mathrm{out}}$
on $(x, y)$ yields the same output as executing $q_{\mathrm{out}}$ on $(x, y)$.
\end{lemma}
